\section{Methods}

\subsection{Wall-Following Control}

The crawler holds a commanded standoff from the pipe wall for the duration of a
traverse. We define the standoff error as
%
\begin{equation}
e(t) = d^{\star} - \hat{d}(t),
\end{equation}
%
where $d^{\star}$ is the commanded standoff and $\hat{d}(t)$ is the estimated distance
from the crawler to the wall. A proportional--derivative law drives this error to zero
through a lateral velocity correction,
%
\begin{equation}
u(t) = \mathrm{sat}\!\left(k_{p}\,e(t) + k_{d}\,\dot{e}(t),\; u_{\max}\right),
\end{equation}
%
where $k_{p}$ and $k_{d}$ are the proportional and derivative gains, $u_{\max}$ bounds
the lateral correction, and
$\mathrm{sat}(x, x_{\max}) = \min(\max(x, -x_{\max}), x_{\max})$. The bound $u_{\max}$
keeps a single correction from stalling one track against the wall. The correction
$u(t)$ is applied as a differential velocity between the two tracks, and the common
track velocity sets the traverse speed. Track velocity commands are issued to the drive
electronics at 100~Hz.

\subsection{Wall-Distance Estimation}

Two time-of-flight sensors mounted on the chassis observe the pipe wall. Their
measurements are low-pass filtered to reject sensor noise and fused by an extended
Kalman filter into the estimate $\hat{d}$ that enters~(1).
% AUTHORS: the draft never states what the filter's state vector contains, nor which
% part of the measurement model is nonlinear (the reason an extended rather than a
% linear filter is required). One sentence on each is needed here. Filter covariances
% and the divergence/iteration guards from the draft are deliberately omitted: they
% are tuning and defensive implementation, and carry no argument.

\subsection{Thermal Compensation}

The time-of-flight sensors drift as their mount heats during a traverse, and this drift
enters $\hat{d}$ as a slowly growing bias. The bias develops on a timescale comparable
to a single traverse, so neither the low-pass filter nor the estimator removes it. We
therefore add a reference sensor at the front bulkhead and correct the two wall-facing
measurements using its output.
% AUTHORS: the compensation law itself is absent from the draft -- what the reference
% sensor observes, and how its reading maps to a correction on the wall-facing
% measurements. This is the paper's contribution and currently has no home in the
% manuscript; it belongs here, as an equation.

This defines the two controllers compared in this work. The \emph{compensated
controller} applies the reference correction. The \emph{uncompensated controller} uses
the wall-facing measurements directly. The two share every other element of the control
software.

\subsection{Experimental Setup}

Experiments were carried out in a 12~m test pipe. The crawler was placed at the pipe
entrance and traversed the pipe autonomously at 0.1~m/s under closed-loop wall
following. Wall distance was logged at the control rate throughout each traverse.

\subsection{Protocol and Analysis}

We performed 16 traverses, alternating between the compensated and the uncompensated
controller, for eight traverses per controller. Performance is reported as the absolute
wall-distance error $|e(t)|$, averaged over a traverse and summarized across traverses
as mean $\pm$ SD. Two traverses in which the error did not reach the 5~mm band before
the midpoint of the pipe were excluded from the aggregate statistics.
% AUTHORS: two points on this exclusion. (1) It reads as post hoc. Either state that
% the criterion was fixed before the campaign, or report all 16 traverses and let the
% two outliers appear in the dispersion -- the latter is the stronger paper.
% (2) State which controller the two excluded traverses ran under; if both were
% uncompensated, the exclusion biases the comparison and a reviewer will find it.
% AUTHORS: the alternating design is paired, so the controller comparison should be
% reported as a paired difference with a 95% CI rather than as two independent means.
% That analysis is not described here because the draft does not report having run it.

\section{Results}

Fig.~3 shows the wall-distance error along a traverse for both controllers. Under the
compensated controller the mean absolute error was 3.1~mm.
% AUTHORS: the aggregate SD is missing from the draft and is required alongside the
% mean. Report as "3.1 $\pm$ X.X mm".

The uncompensated controller held a lower absolute error over the first 30~s of a
traverse. Beyond 30~s the error under the uncompensated controller grew for the
remainder of the traverse, whereas the error under the compensated controller showed no
comparable growth.
% AUTHORS: this paragraph needs the uncompensated controller's aggregate mean $\pm$ SD
% and the difference between controllers expressed as a percentage or ratio. The draft
% reports no aggregate number for the uncompensated controller at all, so the headline
% comparison the paper exists to make is currently unquantified. The directional
% split from the draft (+2.8 $\pm$ 1.1 mm clockwise, -3.4 $\pm$ 1.6 mm
% counterclockwise) is omitted here: signed, direction-split error is an agreement
% analysis, and the claim under test is magnitude of tracking error, not directional
% bias. Restore the split only if direction dependence is itself a finding.

\begin{figure}
\caption{Wall-distance error along a traverse under the compensated and the
uncompensated controller. Curves are the mean across traverses and shaded bands span
one standard deviation. The dashed line marks the 5~mm error band. Values are for eight
traverses per controller.}
\end{figure}
% AUTHORS: the caption decodes curves, bands, and a reference line. Match these to what
% the figure actually plots, and state which color belongs to which controller.

% ---------------------------------------------------------------------------
% MOVED OUT OF RESULTS -- belongs in Discussion. Preserved verbatim in substance so
% nothing is lost; Results carries evidence only, and interpretation is Discussion's
% job.
%
% 1. Interpretation of the compensated result:
%    "The reduction in wall-distance error under the compensated controller supports
%    placing the reference sensor at the front bulkhead, where it observes the same
%    thermal environment as the wall-facing sensors."
%
% 2. Interpretation of the 30 s crossover:
%    "The uncompensated controller degrades only after roughly 30 s, which is
%    consistent with the thermal time constant of the sensor mount."
%    (Hedged to "consistent with" -- one observation matching one mechanism supports
%    possibility, not assertion.)
%
% 3. Limitation, reframed so it states the bound rather than what was unavailable:
%    "Validation was carried out in a clean test pipe. The controller is not restricted
%    to clean pipe, and extending the validation to corroded sections would require
%    only access to a decommissioned line."
% ---------------------------------------------------------------------------
